\documentclass[11pt]{article}
\usepackage[margin=1in]{geometry}
\usepackage{amsmath,amssymb,amsthm}
\usepackage[hidelinks]{hyperref}
\newtheorem{theorem}{Theorem}
\newtheorem{lemma}{Lemma}
\newcommand{\reach}{N_A^2}
\title{Non-symmetry of biconstrained two-step reachability}
\author{Alec Kriebel\\\small Independent researcher\\
\small ORCID: \href{https://orcid.org/0009-0001-9320-500X}{0009-0001-9320-500X}}
\date{October 3, 2026}
\begin{document}
\maketitle
\begin{abstract}
The biconstrained reachability function $\psi(x,y)$ is the largest
fraction of the first part that some vertex of the third part must
reach in two steps, over all finite tripartite graphs satisfying
minimum-degree conditions in both directions on each adjacent pair.
We answer its symmetry question negatively:
$|\psi(13/27,14/27)-\psi(14/27,13/27)|\geq 1/108$.
The proof establishes the exact rational-boundary formula
$\psi(1-\theta,\theta)=1-\theta/d(\theta)$, where $d(\theta)$ is a
minimum integer degree of a positive rational regular incidence
template. An existing seven-type example and its complement bound
both relevant integers by four, making equality impossible.
Neither exact value nor their ordering is needed or asserted.
\end{abstract}

\section{The question}
For $x,y\in(0,1]$, consider finite simple graphs whose vertex set is
partitioned into nonempty stable sets $A,B,C$, with no $A$--$C$ edges,
and with
\begin{align*}
 |N_B(a)|&\geq x|B| &&(a\in A), &
 |N_A(v)|&\geq x|A| &&(v\in B),\\
 |N_C(v)|&\geq y|C| &&(v\in B), &
 |N_B(c)|&\geq y|B| &&(c\in C).
\end{align*}
Such a graph is $(x,y)$-\emph{biconstrained}. Define
\[
 \psi(x,y)=\inf_G\;\max_{c\in C}\frac{|\reach(c)|}{|A|},
 \qquad
 \reach(c)=\{a\in A:N_B(a)\cap N_B(c)\ne\varnothing\},
\]
where the infimum runs over all such finite graphs and partitions.
Reached vertices are counted once, irrespective of the number of paths.
Complete adjacent pairs show that the admissible class is nonempty.
No attainment of this infimum is assumed.

Seymour asked whether $\psi(x,y)=\psi(y,x)$ for all parameters
\cite[p.~47]{S}. The detailed treatment \cite{CHSSS} proves symmetry
for the corresponding one-direction function $\phi$, but explicitly
leaves the biconstrained question unresolved after Theorem~2.3.
Reversing a particular graph exchanges the two directions but does
not by itself equate the universal infima. We prove the following.

\begin{theorem}\label{main}
For the biconstrained function defined above,
\[
 \left|\psi\left(\frac{13}{27},\frac{14}{27}\right)
       -\psi\left(\frac{14}{27},\frac{13}{27}\right)\right|
 \geq\frac{1}{108}.
\]
In particular, $\psi$ is not symmetric.
\end{theorem}

\section{An integer invariant on the rational boundary}
Fix a rational $\theta\in(0,1)$. Let $\mathcal R_\theta$ consist of
finite binary matrices $M$ with nonempty row set $V$ and nonempty
column set $I$, together with strictly positive rational probability
vectors $b=(b_v)_{v\in V}$ and $q=(q_i)_{i\in I}$ satisfying
\begin{equation}\label{regular}
 Mq=\theta\mathbf1_V,\qquad M^{\mathsf T}b=\theta\mathbf1_I.
\end{equation}
Require the columns of $M$ to be pairwise distinct; rows may repeat.
Set
\[
 \Delta(M)=\max_{v\in V}\sum_{i\in I}M_{vi},\qquad
 d(\theta)=\min_{(M,b,q)\in\mathcal R_\theta}\Delta(M).
\]
The degree $\Delta$ counts distinct column types without weights.
If $\theta=p/r$ in lowest terms, the $r$ cyclic intervals of length
$p$ on $r$ positions, with both weights uniform $1/r$, give distinct
columns and degree $p$. Thus $\mathcal R_\theta$ is nonempty and
$1\leq d(\theta)\leq p$. By well-ordering of positive integers,
this minimum is realized by a finite template. This is a separate
fact from attainment of the graph infimum defining $\psi$.

\begin{theorem}\label{boundary}
For every rational $\theta\in(0,1)$,
\begin{equation}\label{formula}
 \psi(1-\theta,\theta)=1-\frac{\theta}{d(\theta)}.
\end{equation}
The value on this boundary is attained by a finite graph.
\end{theorem}

\begin{proof}
First take any finite $(1-\theta,\theta)$-biconstrained graph. Put
\[
 m=|B|,\qquad F=\max_{c\in C}\frac{|\reach(c)|}{|A|}.
\]
If $F=1$, the required lower bound is immediate.
Suppose $F<1$. Every $c\in C$ misses
some $a\in A$. Their $B$-neighborhoods are disjoint and have sizes
at least $\theta m$ and $(1-\theta)m$. Hence they partition $B$,
with both size bounds attained. In particular,
$|N_B(c)|=\theta m$ for every $c$.
Every row has at least $\theta|C|$ neighbors in $C$.
The $B$--$C$ edge count is exactly $\theta m|C|$,
so every row has exactly $\theta|C|$ neighbors.

Group the vertices of $C$ by their distinct $B$-neighborhoods
$T_1,\ldots,T_n$. Let $M_{vi}=\mathbf1_{v\in T_i}$, let
$b_v=1/m$, and let $q_i$ be the fraction of $C$ in class $i$.
These positive rational weights satisfy \eqref{regular}.
For each $i$, the vertices missed by a $c$ of class $i$ are exactly
\[
 A_i=\{a\in A:N_B(a)=B\setminus T_i\}.
\]
Indeed, disjointness from $T_i$ and the minimum degree force this
equality of neighborhoods. The $A_i$ are distinct, disjoint and
nonempty. Writing $p_i=|A_i|/|A|$ and $t=\min_i p_i$ gives
$F=1-t$. For every $v\in B$, all $A_i$ with $v\in T_i$ are
nonneighbors of $v$, so the $B$-to-$A$ constraint implies
\[
 \sum_{i:M_{vi}=1}p_i\leq\theta,
 \qquad t\Delta(M)\leq\theta.
\]
Additional vertices of $A$ outside these classes are retained;
they cannot weaken this inequality. As $\Delta(M)\geq d(\theta)$,
$F\geq1-\theta/d(\theta)$ for every admissible graph.

For the upper bound choose a template realizing $d=d(\theta)$,
with $n$ columns and supports $T_i\subseteq V$. Every support
has $b$-weight $\theta$. Strict positivity of $b$ and distinctness
of the columns make the supports pairwise incomparable.
Set $t=\theta/d$ and $\delta=1-nt$. The residual is nonnegative,
since
\[
 n\theta=\sum_i b(T_i)=\sum_v b_v\sum_i M_{vi}\leq d.
\]
Build a weighted graph with row vertices $v$ of weight $b_v$ in $B$,
vertices $c_i$ of weight $q_i$ in $C$ adjacent to $T_i$, and vertices
$a_i$ of weight $t$ in $A$ adjacent to $V\setminus T_i$.
If $\delta>0$, add a vertex $a_*$ of weight $\delta$ adjacent
to all of $B$; if $\delta=0$, omit it. Each part has weight one
and every retained type has positive weight.

Each $a_i$ sees $B$-weight $1-\theta$, while $a_*$ sees one.
Each row $v$ sees $A$-weight
$1-t\sum_iM_{vi}\geq1-\theta$. Both directions between $B$
and $C$ have neighbor weight exactly $\theta$ by \eqref{regular}.
Thus all four constraints hold. The vertex $c_i$ misses $a_i$,
but reaches every $a_j$ for $j\ne i$, since
$T_i\setminus T_j\ne\varnothing$. It also reaches $a_*$ when
present, because $\theta>0$. Its reached $A$-weight is therefore
exactly $1-t$.

Finally, replace each weighted type by a positive integer number
of twins proportional to its rational weight, using a common
denominator within each part. Join entire type pairs according to
the weighted adjacency. All middle types survive, so both the
degree fractions and the reachability fractions are preserved
exactly. This is an ordinary finite simple graph attaining
$F=1-t$, which proves \eqref{formula} and attainment.
\end{proof}

\section{The asymmetric pair}
We use the regular bipartite template from Figure~1 of
\cite[p.~6; also Section~12, p.~71]{CHSSS}. Its matrix and weights are
\[
 M=\begin{pmatrix}
 1&0&1&1&1&0&0\\
 0&1&1&1&1&0&0\\
 0&0&1&0&0&1&1\\
 0&0&0&1&0&1&1\\
 0&0&0&0&1&1&1\\
 1&1&0&0&0&1&0\\
 1&1&0&0&0&0&1
 \end{pmatrix},\qquad
 b=\frac{(5,5,3,3,3,4,4)^{\mathsf T}}{27},\quad
 q=\frac{(4,4,3,3,3,5,5)^{\mathsf T}}{27}.
\]
Direct multiplication gives
$Mq=(13/27)\mathbf1$ and $M^{\mathsf T}b=(13/27)\mathbf1$.
The columns are distinct and the maximum row degree is four.
The complement $J-M$, where $J$ is the all-ones matrix, has the
same positive weights, regularity $14/27$, distinct columns,
and maximum row degree four. Hence, with
$d=d(13/27)$ and $e=d(14/27)$,
\[
 d,e\in\{1,2,3,4\},\qquad
 \psi(14/27,13/27)=1-\frac{13}{27d},\quad
 \psi(13/27,14/27)=1-\frac{14}{27e}.
\]
Equality would require $13e=14d$, impossible in this range.
For the quantitative gap, if $d=e$ the difference is
$1/(27d)\geq1/108$. If $e>d$, then
$13e-14d=13(e-d)-d\geq10$ and $27de\leq324$.
If $d>e$, then $14d-13e=14(d-e)+e\geq15$ with the same
denominator bound. Both latter cases give a larger gap.
This proves Theorem~\ref{main}.

\paragraph{Explicit examples and their limits.}
The construction above applied to $M$ with degree four gives
part sizes $(108,27,27)$, $A$-type sizes $(13,13,13,13,13,13,13,17)$,
and $B,C$-type sizes given by the numerators of $b,q$.
Every $c$ reaches $95$ vertices of $A$.
The complement gives $A$-type sizes $(14,14,14,14,14,14,14,10)$
and reaches $94$ vertices of $A$.
Thus $\psi(14/27,13/27)\leq95/108$ and
$\psi(13/27,14/27)\leq47/54$.
These are upper bounds, not exact evaluations; their inequality
alone would not prove asymmetry. The universal boundary formula
is indispensable. The true integers $d,e$, the exact values,
and their ordering remain unevaluated here.

\section{Prior work, verification, and disclosure}
The finite weighted-graph framework and the seven-type matrix
are prior work \cite{CHSSS}, not new constructions of this note.
Theorem~12.2 there concerns symmetry when $\psi(x,y)$ attains
its minimal value $x$; it does not establish general symmetry.
The contribution here is the integer boundary formula, with
universal lower and finite upper proofs, and its deduction of
non-symmetry at an explicit rational pair.
A dated literature audit accompanies the supplement. Its search
has finite coverage and does not certify absence of all prior work.

The supplement contains the complete original proof records,
exact arithmetic checkers, two independent mathematical audits,
source identities and a portable verifier. It checks the matrix,
both ordinary graphs, all sixteen possible integer gaps and
controls for invalid inferences. The theorem is an elementary
argument over all finite graphs; finite tests and file
hashes supplement rather than replace its proof. No external
optimization theorem, numerical approximation, or presumed
extremizer for $\psi$ is required.

\paragraph{AI use and review status.}
AI tools were used extensively to develop and verify the proofs,
write and run checking code, search the literature, draft this note,
and conduct adversarial reviews. Independent AI review paths
are documented in the accompanying research records.
This manuscript is an unrefereed preprint and has not received
external human peer review.

\begin{thebibliography}{9}
\raggedright
\bibitem{S}
P. Seymour, \emph{Concatenating bipartite graphs},
Oberwolfach Reports 1/2019, pp.~46--47.
\href{https://doi.org/10.4171/OWR/2019/1}{doi:10.4171/OWR/2019/1}.
\bibitem{CHSSS}
M. Chudnovsky, P. Hompe, A. Scott, P. Seymour, and S. Spirkl,
\emph{Concatenating Bipartite Graphs},
Electronic Journal of Combinatorics 29(2) (2022), \#P2.47.
\href{https://doi.org/10.37236/8451}{doi:10.37236/8451}.
\end{thebibliography}
\end{document}
